% 缩略词表

\newpage
\begingroup
  \ctexset
    {
      chapter =
        {
          format += \centering,
          numbering = false,
          pagestyle = bachelor-frontmatter
        }
    }
  \chapter{缩略词表}
\endgroup

\begingroup
\zihao{5}
\setlength{\tabcolsep}{3pt}
\setlength{\LTpre}{0pt}
\setlength{\LTpost}{0pt}
\renewcommand{\arraystretch}{1.08}
\begin{longtable}
  {
    |>{\centering\arraybackslash}m{1.95cm}
    |>{\centering\arraybackslash}m{4.70cm}
    |>{\centering\arraybackslash}m{1.95cm}
    |>{\centering\arraybackslash}m{4.70cm}|
  }
  \hline
  \textbf{缩略词} & \textbf{含义或本文用法} &
  \textbf{缩略词} & \textbf{含义或本文用法} \\
  \hline
  \endfirsthead
  \hline
  \textbf{缩略词} & \textbf{含义或本文用法} &
  \textbf{缩略词} & \textbf{含义或本文用法} \\
  \hline
  \endhead
  \hline
  \endfoot
  \hline
  \endlastfoot
  AI & 人工智能；本文主要指神经网络前处理方法 &
  AR/VR & 增强现实/虚拟现实应用场景 \\
  \hline
  AV1 & AOMedia Video 1开放视频编码格式 &
  AVC & 高级视频编码，通常对应H.264标准 \\
  \hline
  AVS & 音视频编码标准体系 &
  BasicVSR++ & 高性能视频复原模型，本文作为复杂度参照 \\
  \hline
  BD-Rate & Bjontegaard平均码率差异指标 &
  BD-LPIPS & 以LPIPS为质量轴的BD码率差异指标 \\
  \hline
  CNN & 卷积神经网络 &
  CQP & 恒定量化参数码率控制模式 \\
  \hline
  CTC & 通用测试条件；本文指HEVC测试序列设置 &
  CTU & 编码树单元，HEVC中的基本块划分单元 \\
  \hline
  CUDA & NVIDIA通用并行计算平台 &
  CUVID & NVIDIA硬件视频解码接口 \\
  \hline
  DBF & 去块效应滤波器 &
  DCT & 离散余弦变换 \\
  \hline
  DDPG & 深度确定性策略梯度算法 &
  DQN & 深度Q网络算法 \\
  \hline
  EA-MCTF & 编码器感知运动补偿时域滤波方法 &
  FFmpeg & 通用音视频处理框架 \\
  \hline
  GAE & 广义优势估计 &
  GOP & 图像组，表示关键帧间隔及预测结构 \\
  \hline
  GPU & 图形处理器；本文用于模型推理和硬件编解码 &
  HEVC & 高效率视频编码，通常对应H.265标准 \\
  \hline
  HVS & 人眼视觉系统 &
  LPIPS & 学习感知图像块相似度指标 \\
  \hline
  MACs & 乘加运算量，模型计算复杂度指标 &
  MFQE & 多帧质量增强方法 \\
  \hline
  MSE & 均方误差 &
  NPP & NVIDIA性能原语库 \\
  \hline
  NVDEC & NVIDIA硬件视频解码器 &
  NVENC & NVIDIA硬件视频编码器 \\
  \hline
  PSNR & 峰值信噪比 &
  QP & 量化参数 \\
  \hline
  RAPNet & 码率自适应前处理网络，本文Actor骨干网络 &
  RD & 率失真关系或率失真性能 \\
  \hline
  RL & 强化学习 &
  RPP & 率感知预处理方法 \\
  \hline
  RVRT & 循环视频复原Transformer &
  SAC & 软Actor-Critic算法 \\
  \hline
  SAO & 样本自适应偏移滤波器 &
  SDK & 软件开发工具包 \\
  \hline
  SSIM & 结构相似性指标 &
  TV & 全变分，本文用于平滑正则项 \\
  \hline
  VMAF & 视频多方法评价融合指标 &
  VRT & 视频复原Transformer \\
  \hline
  VVC & 通用视频编码，通常对应H.266标准 &
  YCbCr & 亮度、蓝色色度差与红色色度差颜色表示 \\
  \hline
  YUV & 亮度与色度分量表示 &
  YUV420 & YUV 4:2:0色度二分之一采样格式 \\
\end{longtable}
\endgroup
\newpage
